\documentclass{article}
\usepackage[margin=1in]{geometry}
\usepackage{graphicx}
\usepackage{booktabs}
\usepackage{float}
\usepackage{amsmath}
\usepackage{hyperref}

\title{IFT6135 - Assignment 2 Report}
\author{Guillaume Genois, 20248507}
\date{March 20, 2026}

\begin{document}

\maketitle

%══════════════════════════════════════════════════════════════════════════════
\section*{Problem 2: Training Sequential Models}
%══════════════════════════════════════════════════════════════════════════════

% ─────────────────────────────────────────────────────────────────────────────
\subsection*{Question 1: Learning curves (6 pts)}

For each of the 8 experiments, the training loss and validation loss/accuracy
are plotted over training iterations.

% TODO: Add 8 learning curve figures (one per experiment or grouped).
% Example figure template:
%
% \begin{figure}[H]
%   \centering
%   \includegraphics[width=\linewidth]{exp1_curves.png}
%   \caption{Experiment \textbf{1}: Training loss (left) and validation
%            loss/accuracy (right) over training iterations.}
%   \label{fig:exp1_curves}
% \end{figure}

\textit{TODO: Insert learning curve plots for experiments 1--8.}

% ─────────────────────────────────────────────────────────────────────────────
\subsection*{Question 2: Summary table (5 pts)}

\begin{table}[H]
  \centering
  \caption{Summary of train and validation performance for all 8 experiments.
           Best validation result is shown in \textbf{bold}. ``Train time''
           is the total wall-clock training time; ``Eval time'' is the time
           to evaluate the model on the validation set.}
  \label{tab:summary}
  \begin{tabular}{clrrrrr}
    \toprule
    \textbf{Exp.} & Architecture & Train Acc & Val Acc & Train Loss & Train Time & Eval Time \\
    \midrule
    \textbf{1} & TODO & -- & -- & -- & -- & -- \\
    \textbf{2} & TODO & -- & -- & -- & -- & -- \\
    \textbf{3} & TODO & -- & -- & -- & -- & -- \\
    \textbf{4} & TODO & -- & -- & -- & -- & -- \\
    \textbf{5} & TODO & -- & -- & -- & -- & -- \\
    \textbf{6} & TODO & -- & -- & -- & -- & -- \\
    \textbf{7} & TODO & -- & -- & -- & -- & -- \\
    \textbf{8} & TODO & -- & -- & -- & -- & -- \\
    \bottomrule
  \end{tabular}
\end{table}

% ─────────────────────────────────────────────────────────────────────────────
\subsection*{Question 3: Wall-clock time vs.\ generalization (2 pts)}

\textit{TODO: Among the 8 configurations, which would you use if you were most
concerned with wall-clock time? With generalization performance? Justify.}

% ─────────────────────────────────────────────────────────────────────────────
\subsection*{Question 4: Impact of dropout (3 pts)}

\textit{TODO: Compare experiments \textbf{1} and \textbf{2}. What was the
impact of having dropout?}

% ─────────────────────────────────────────────────────────────────────────────
\subsection*{Question 5: Encoder vs.\ Encoder-Decoder (3 pts)}

\textit{TODO: Compare experiments \textbf{2} and \textbf{3}. Are there
benefits of having an Encoder-Decoder network instead of just an Encoder?
Speculate as to why or why not.}

% ─────────────────────────────────────────────────────────────────────────────
\subsection*{Question 6: Impact of soft-attention (3 pts)}

\textit{TODO: Compare experiments \textbf{3} and \textbf{4}. What was the
impact of training the GRU network with soft-attention?}

% ─────────────────────────────────────────────────────────────────────────────
\subsection*{Question 7: Transformer experiments (3 pts)}

\textit{TODO: In experiments \textbf{5}, \textbf{6}, \textbf{7} and
\textbf{8}, you trained or finetuned a Transformer. Given the recent high
profile Transformer-based models, are the results as you expected? Speculate
as to why or why not.}

% ─────────────────────────────────────────────────────────────────────────────
\subsection*{Question 8: GPU memory usage (3 pts)}

% TODO: Add GPU memory table
%
% \begin{table}[H]
%   \centering
%   \caption{Average steady-state GPU memory usage for each experiment.}
%   \label{tab:gpu_memory}
%   \begin{tabular}{clr}
%     \toprule
%     \textbf{Exp.} & Architecture & GPU Memory (MB) \\
%     \midrule
%     \textbf{1} & TODO & -- \\
%     ...
%     \bottomrule
%   \end{tabular}
% \end{table}

\textit{TODO: For each experiment configuration, measure the average
steady-state GPU memory usage. Comment about the GPU memory footprints of each
model, discussing reasons behind increased or decreased memory consumption
where applicable.}

% ─────────────────────────────────────────────────────────────────────────────
\subsection*{Question 9: Learning curve analysis (7 pts)}

\textit{TODO: Comment on the learning curve behavior of the various models
trained, under different hyper-parameter settings. Did a particular class of
models over/underfit more easily than the others? Exhibited instabilities in
training? Can you make an informed guess of the various steps a practitioner
can take to prevent over/under-fitting or instability in this case?}

% ─────────────────────────────────────────────────────────────────────────────
\subsection*{Question 10: Cross Attention vs.\ Self Attention (5 pts)}

\textit{TODO: Give a high-level overview of what an attention mechanism is,
and go into each attention mechanism (Self, Cross), and their differences.
Write two paragraphs total.}

\end{document}
